% sections/02-sqlmap-activa.tex
\section{Inyección SQL activa con sqlmap}

\subsection{Selección del objetivo}

Primeramente se realizó una búsqueda avanzada en Google con el \eng{dork} \texttt{inurl: .asp?id}, orientado a identificar páginas con parámetros potencialmente inyectables; para el ejercicio puede hacerse uso de una URL propia identificada en la búsqueda o de cualquiera de las que propone la guía del laboratorio:

\begin{itemize}
  \item \url{http://testphp.vulnweb.com/listproducts.php?cat=1}
  \item \url{https://www.china-mould.com/cn_Contact.asp?ID=12}
\end{itemize}

Sobre esta última se trabajó la práctica, accediendo en Kali a la categoría \texttt{APLICACIONES/04} (evaluación de bases de datos) y, estando en la línea de comando, a la utilidad sqlmap \cite{sqlmap-docs}.

\subsection{Opciones de la herramienta}

Estando en la línea de comando se ejecutó la opción de ayuda:

\begin{lstlisting}[language=bash]
sqlmap -h
\end{lstlisting}

\pregunta{Describa las opciones que pueden ser usadas por esta herramienta.}
\begin{respuesta}
  De la salida de la ayuda se destacan las opciones que se usaron (y otras relevantes) a lo largo del ejercicio:

  \begin{center}
  \small
  \begin{tabularx}{\linewidth}{>{\ttfamily}l X}
    \toprule
    \normalfont\bfseries Opción & \bfseries Descripción \\
    \midrule
    -u URL & Objetivo de la evaluación, con el parámetro que se sospecha inyectable incluido en la URL. \\
    --banner & Recupera el banner del DBMS (versión del motor de base de datos). \\
    --current-user & Recupera el usuario actual con el que la aplicación se conecta a la base de datos. \\
    --current-db & Recupera la base de datos actual (instancia en uso). \\
    --users & Enumera los usuarios del gestor de base de datos. \\
    --passwords & Enumera y recupera los hashes de las contraseñas de los usuarios del gestor. \\
    --dbs & Enumera las bases de datos disponibles. \\
    -D BD & Delimita la base de datos objetivo de las siguientes opciones. \\
    --tables & Enumera las tablas (junto con \texttt{-D}). \\
    -T tabla & Delimita la tabla objetivo de las siguientes opciones. \\
    --columns & Enumera las columnas de una tabla (con \texttt{-D} y \texttt{-T}). \\
    --dump & Volca el contenido de una tabla (con \texttt{-D} y \texttt{-T}). \\
    --batch & Ejecuta sin preguntar, tomando el comportamiento por defecto en cada decisión interactiva. \\
    --level / --risk & Profundidad y agresividad de las pruebas (número de parámetros y payloads probados). \\
    -p parámetro & Delimita el parámetro a testear, en lugar de probarlos todos. \\
    --forms & Analiza los formularios de la página objetivo en busca de campos inyectables. \\
    --cookie / --headers & Permite autenticar la sesión (cookies de sesión) o cabeceras personalizadas. \\
    --dbms & Fuerza el motor de base de datos, evitando el fingerprinting inicial. \\
    --os-shell & Intenta obtener una consola interactiva del sistema operativo del servidor. \\
    \bottomrule
  \end{tabularx}
  \end{center}
\end{respuesta}

\porque{antes de continuar con la tabla}{%
  La tabla describe opciones documentadas de sqlmap y es verificable contra tu
  propia salida de \texttt{sqlmap -h}; revísala ahí, quita las que no aparezcan
  o añade las que te llamen la atención, y sobre todo deja claro cuáles usaste
  en los ítems siguientes. Ojo con el encoding: copia la tabla tal cual, ya
  está escrita para no romper el ancho de columna.%
}

\subsection{Escaneo inicial del objetivo}

Procediendo con el escaneo de la URL identificada, se ejecutó la sentencia que propone la guía:

\begin{lstlisting}[language=bash]
sqlmap -u "http://testphp.vulnweb.com/listproducts.php?cat=1" --batch
\end{lstlisting}

\captura{sqlmap-escaneo.png}{Salida del escaneo inicial: parámetro inyectable, servidor web y gestor de base de datos identificados}

\begin{analisis}
\end{analisis}

\porque{qué debe verse en la captura}{%
  En caso de encontrar parámetros inyectables, sqlmap despliega al final un
  resumen con el \emph{back-end DBMS} (se espera MySQL para
  testphp.vulnweb.com), el servidor web (usualmente nginx en ese host de
  Acunetix) y el tipo de inyección encontrada. Base para el análisis: ``La
  salida evidencia que el parámetro \texttt{cat} es inyectable, confirmando que
  la aplicación concatena directamente el valor recibido sobre la consulta SQL
  sin ningún tipo de sanitización, de forma que un atacante puede alterar la
  semántica de la sentencia que se ejecuta del lado del motor.'' Completa con
  los valores exactos de tu ejecución.%
}

\subsection{Versión del motor de la base de datos}

\begin{lstlisting}[language=bash]
sqlmap -u "http://testphp.vulnweb.com/listproducts.php?cat=1" --banner --batch
\end{lstlisting}

\captura{sqlmap-banner.png}{Banner del DBMS: versión del motor de la base de datos}

\begin{analisis}
\end{analisis}

\porque{guía para este ítem}{%
  Con \texttt{--banner} sqlmap fuerza la recuperación del banner del motor
  (equivalente a un \texttt{SELECT @@VERSION} en MySQL), el cual se espera que
  informe MySQL como DBMS con su versión exacta. Base: ``El banner obtenido
  revela el motor y la versión exacta que ejecuta el servidor, información que
  resulta valiosa para el atacante porque le permite seleccionar payloads y
  técnicas de explotación específicas para esa versión ( exploits públicos,
  errores conocidos), adelantando incluso la fase de post-explotación.''%
}

\subsection{Usuario actual de la base de datos}

\begin{lstlisting}[language=bash]
sqlmap -u "http://testphp.vulnweb.com/listproducts.php?cat=1" --current-user --batch
\end{lstlisting}

\captura{sqlmap-usuario.png}{Usuario actual con el que la aplicación se conecta al gestor}

\begin{analisis}
\end{analisis}

\porque{guía para este ítem}{%
  Base: ``El usuario actual corresponde a la cuenta con la que la aplicación
  web se autentica contra el motor de base de datos (para este objetivo suele
  verse el usuario \texttt{ACUART}); su relevancia radica en que, de resultar
  un usuario con privilegios administrativos (DBA), el atacante obtendría
  acceso a operaciones sobre todo el gestor y no solo sobre el esquema de la
  aplicación.'' Verifica el nombre exacto en tu salida.%
}

\subsection{Base de datos actual (instancia)}

\begin{lstlisting}[language=bash]
sqlmap -u "http://testphp.vulnweb.com/listproducts.php?cat=1" --current-db --batch
\end{lstlisting}

\captura{sqlmap-instancia.png}{Instancia (base de datos) en uso por la aplicación}

\begin{analisis}
\end{analisis}

\porque{guía para este ítem}{%
  Se espera que la instancia actual sea \texttt{acuart}; con ese nombre se
  delimitan los comandos siguientes (\texttt{-D acuart}). Base: ``Identificar
  la instancia actual permite al atacante focalizar la enumeración sobre el
  esquema real de la aplicación, encadenando posteriormente el listado de sus
  tablas y columnas sin necesidad de recorrer todas las bases de datos del
  servidor.''%
}

\subsection{Hashes de las contraseñas de los usuarios}

\begin{lstlisting}[language=bash]
sqlmap -u "http://testphp.vulnweb.com/listproducts.php?cat=1" --passwords --batch
\end{lstlisting}

\captura{sqlmap-hashes.png}{Usuarios del gestor y hashes de sus contraseñas recuperados}

\begin{analisis}
\end{analisis}

\porque{guía para este ítem}{%
  \texttt{--passwords} enumera los usuarios del gestor e intenta leer sus
  hashes; lo que devuelva depende de los privilegios del usuario actual (si
  no es DBA puede fallar o devolver solo el usuario actual, lo cual también es
  un hallazgo válido y hay que reportarlo). Base: ``La recuperación de hashes
  de contraseñas constituye una de las afectaciones más serias evidenciadas,
  dado que estos pueden someterse a cracking fuera de línea (diccionario o
  fuerza bruta) y, de reutilizarse la misma credencial en otros sistemas,
  escalar el incidente más allá de la aplicación vulnerable.'' Si sqlmap
  pregunta por almacenar la tabla diccionario, acepta con el default de
  \texttt{--batch}.%
}

\subsection{Esquema de tablas de la base de datos actual}

\begin{lstlisting}[language=bash]
sqlmap -u "http://testphp.vulnweb.com/listproducts.php?cat=1" -D acuart --tables --batch
\end{lstlisting}

\captura{sqlmap-esquema.png}{Tablas de la instancia acuart enumeradas}

\begin{analisis}
\end{analisis}

\porque{guía para este ítem}{%
  Ajusta \texttt{-D} al nombre real que te devolvió el ítem de la instancia.
  Base: ``La enumeración del esquema expone la estructura interna de la
  aplicación (se esperan tablas como \texttt{artists}, \texttt{carts} o
  \texttt{customers}), información que en manos del atacante define
  exactamente dónde resides la información sensible y habilita el volcado
  dirigido de cualquiera de ellas.''%
}

\subsection{Columnas de una de las tablas}

\begin{lstlisting}[language=bash]
sqlmap -u "http://testphp.vulnweb.com/listproducts.php?cat=1" -D acuart -T artists --columns --batch
\end{lstlisting}

\captura{sqlmap-columnas.png}{Columnas y tipos de la tabla artists}

\begin{analisis}
\end{analisis}

\porque{guía para este ítem}{%
  Puedes cambiar \texttt{artists} por la tabla que te llame la atención
  (\texttt{customers} suele ser más ilustrativa por contener datos personales).
  Base: ``El detalle de columnas y tipos de dato completa el mapa del esquema,
  permitiendo al atacante construir consultas quirúrgicas (por ejemplo,
  seleccionando únicamente las columnas de credenciales o de datos personales)
  y reduciendo el ruido de la extracción para evitar detección.''%
}

\subsection{Volcado del contenido de una tabla}

\begin{lstlisting}[language=bash]
sqlmap -u "http://testphp.vulnweb.com/listproducts.php?cat=1" -D acuart -T artists --dump --batch
\end{lstlisting}

\captura{sqlmap-dump.png}{Registros de la tabla artists volcados en pantalla}

\begin{analisis}
\end{analisis}

\porque{guía para este ítem}{%
  Base: ``El volcado materializa la exfiltración: los registros salen del
  motor y quedan en poder del atacante (sqlmap incluso los persiste en un CSV
  local con \texttt{--dump}); sobre una tabla con datos reales de clientes el
  mismo comando entregaría nombres, correos y tarjetas, con lo cual queda
  demostrado que la falta de parametrización de una sola consulta compromete
  la confidencialidad de toda la información hospedada.'' De haber volcado
  \texttt{customers}, cierra la cadena con la mención de qué datos personales
  quedaron expuestos.%
}
